\documentclass{article}

\usepackage{amsmath,amssymb,amsthm}
\usepackage{booktabs}
\usepackage{multirow}
\usepackage{graphicx}
\usepackage{hyperref}
\usepackage{microtype}
\usepackage{xcolor}
\usepackage{cleveref}

\title{Analytic Surrogates for Pathwise Robustness:\\
On the Promise and Limits of Rouch\'e-Type Zero-Counting\\
in Neural Network Verification}
\author{Author Names Omitted}
\date{}

\newtheorem{theorem}{Theorem}
\newtheorem{proposition}{Proposition}
\newtheorem{lemma}{Lemma}
\newtheorem{corollary}{Corollary}
\newtheorem{remark}{Remark}
\theoremstyle{definition}
\newtheorem{definition}{Definition}

\newcommand{\C}{\mathbb{C}}
\newcommand{\R}{\mathbb{R}}
\newcommand{\norm}[1]{\left\|#1\right\|}
\newcommand{\abs}[1]{\left|#1\right|}
\DeclareMathOperator*{\argmax}{arg\,max}
\DeclareMathOperator*{\argmin}{arg\,min}

\newcommand{\TBD}[1]{\textcolor{red}{[TBD: #1]}}

\begin{document}
\maketitle

\begin{abstract}
We investigate whether complex-analytic zero-counting can soundly certify
the absence of classifier boundary crossings along input paths.  Given a
continuous path $\gamma:[0,1]\to\mathbb{R}^d$ between two same-class inputs,
the per-class margin $g_j(t)$ must remain positive throughout.  We construct
polynomial surrogates $Q_j$ by Chebyshev sampling, bound the approximation
error $\epsilon_j$, and apply Rouch\'e's theorem on a rectangular complex
domain to certify zero-freeness.  On synthetic polynomial classifiers with
controlled margin and degree, the real-polynomial certificate achieves 100\%
across all settings (exact surrogates, zero approximation error), while the
Rouch\'e certificate achieves only 0--46\% even with exact surrogates: the
failure is geometric, not error-driven, as higher-degree polynomial surrogates
take large complex values off the real axis that defeat the constant-reference
Rouch\'e condition.  On a small MLP
trained on the two-moons dataset, certification rates collapse rapidly as
path length is scaled beyond the original interpolation, because the
surrogate approximation error grows faster than the available margin.
Our main finding is a sharp negative result: Rouch\'e certificates offer no
advantage over the real-polynomial check on real-analytic paths, and both
become vacuous on realistic neural paths beyond short interpolation scales.
We provide a counterexample showing why naive complexification fails, a sound
theorem connecting zero-freeness to pathwise positivity, and practical
guidelines for when complex-analytic certification is not worth attempting.
\end{abstract}

%% ============================================================
\section{Introduction}
\label{sec:intro}

Post-hoc certification of neural network behaviour is an active area with
two broad strategies: \emph{bound propagation} (interval arithmetic, abstract
interpretation) and \emph{sampling with verification} (Lipschitz bounds,
Lyapunov methods).  A less-explored direction asks whether classical
complex analysis---specifically Rouch\'e's theorem on zero-counting---can
provide certificates for the absence of class-boundary crossings along
continuous input paths.

The idea is appealing: fit an analytic surrogate $Q_j(t)$ to the margin
function $g_j(t) = f_c(\gamma(t)) - f_j(\gamma(t))$ on $[0,1]$, extend
analytically to a complex domain $D\supset[0,1]$, and apply Rouch\'e to
certify zero-freeness in $D$.  Since $[0,1]\subset D$, zero-freeness in $D$
implies zero-freeness on $[0,1]$, and one sign evaluation completes the
positivity certificate.

Papers~1 and~2 of this series developed Rouch\'e-type certificates for
the spectral radius of learned transition matrices---a setting where the
``path'' is the unit circle and the complex structure is intrinsic.  Paper~3
asks whether the same toolkit transfers to the classifier boundary problem,
where the analytic structure must be \emph{imposed} through a surrogate.

Our contributions are:
\begin{enumerate}
  \item A sound theorem (Theorem~\ref{thm:main}) connecting Rouch\'e
    zero-freeness to pathwise margin positivity via a surrogate error bound.
  \item A counterexample (Section~\ref{sec:counterexample}) showing that naive
    analytic continuation can introduce spurious zeros not present on the real path.
  \item Experiments on synthetic polynomial classifiers
    (Section~\ref{sec:exp_synthetic}) showing the real-polynomial certificate
    achieves 100\% with exact surrogates while Rouché fails even then (0--46\%),
    establishing that the Rouché failure is geometric rather than error-driven.
  \item Experiments on a small MLP (Section~\ref{sec:exp_networks}) showing
    certification collapses with path length, and that Rouch\'e offers no
    improvement over the simpler real-polynomial check.
  \item Practical guidelines (Section~\ref{sec:discussion}) for when
    complex-analytic path certification is not worth attempting.
\end{enumerate}

%% ============================================================
\section{Background and Setup}
\label{sec:setup}

\subsection{Classifier path margin}

Let $f:\mathbb{R}^d\to\mathbb{R}^K$ be a classifier, $x_a, x_b$ inputs
predicted as class $c$, and $\gamma:[0,1]\to\mathbb{R}^d$ a continuous path
with $\gamma(0)=x_a$, $\gamma(1)=x_b$.  For each rival class $j\neq c$ define
\begin{equation}
  g_j(t) = f_c(\gamma(t)) - f_j(\gamma(t)).
  \label{eq:margin}
\end{equation}
The path is \emph{class-preserving} if $g_j(t)>0$ for all $t\in[0,1]$ and
all $j\neq c$.

\subsection{Rouch\'e's theorem}

\begin{theorem}[Rouch\'e, classical]
  Let $D\subset\mathbb{C}$ be a simply connected domain with boundary
  $\partial D$.  If $F,G$ are analytic on $\bar D$ and
  $\abs{G(z)}<\abs{F(z)}$ for all $z\in\partial D$, then $F$ and $F+G$
  have the same number of zeros (counted with multiplicity) in $D$.
\end{theorem}

We apply this with $F=R_j$ (a reference analytic function with no zeros in
$D$) and $G=Q_j-R_j$, concluding that $Q_j$ has no zeros in $D$ whenever
$\abs{R_j(z)}>\abs{Q_j(z)-R_j(z)}$ on $\partial D$.

%% ============================================================
\section{Analytic Surrogate Construction}
\label{sec:surrogate}

\subsection{Chebyshev polynomial surrogate}

Let $g_j:[0,1]\to\mathbb{R}$ be the margin function.  We fit a degree-$d$
polynomial surrogate $Q_j$ by evaluating $g_j$ at the $n\geq d+1$ Chebyshev
nodes of the first kind
\begin{equation}
  t_k = \frac{1}{2}\!\left(1 + \cos\frac{(2k-1)\pi}{2n}\right),
  \quad k=1,\ldots,n,
\end{equation}
and solving the least-squares system.  Chebyshev sampling minimises Runge
oscillation and gives near-optimal approximation rates for smooth $g_j$.

\subsection{Approximation error bound}

After fitting $Q_j$, we estimate
\begin{equation}
  \epsilon_j \;=\; \max_{t\in\mathcal{T}} \abs{g_j(t) - Q_j(t)},
\end{equation}
where $\mathcal{T}$ is a uniform grid of 500 points on $[0,1]$.  This is a
numerical estimate, not a rigorous bound; obtaining a rigorous bound would
require Bernstein ellipse analysis or interval arithmetic.  We treat
$\epsilon_j$ as a conservative estimate and note this caveat throughout.

%% ============================================================
\section{Certificates}
\label{sec:certs}

\subsection{Real-polynomial positivity certificate}

A direct real certificate requires no complex analysis:
\begin{equation}
  \min_{t\in[0,1]} Q_j(t) > \epsilon_j
  \;\Longrightarrow\; g_j(t) > 0 \;\text{for all}\; t\in[0,1].
\end{equation}
This is tight: the condition fails iff the surrogate minimum is within the
approximation error of zero.

\subsection{Rouch\'e zero-free certificate}

Let $D\subset\mathbb{C}$ be a rectangular domain enclosing $[0,1]$:
\begin{equation}
  D = \{z\in\mathbb{C} : -w \leq \operatorname{Re}(z) \leq 1+w,\;
        \abs{\operatorname{Im}(z)} \leq w\}
\end{equation}
for half-width $w>0$.  Let $c\in\mathbb{C}$ be the centroid of $Q_j$ on
$\partial D$, used as a constant zero-free reference $R_j\equiv c$.

\begin{definition}[Rouch\'e margin]
  \begin{equation}
    m_{\mathrm{R}}(Q_j,\partial D) \;=\; \abs{c} - \max_{z\in\partial D}\abs{Q_j(z)-c}.
    \label{eq:rouche_margin}
  \end{equation}
\end{definition}

\begin{theorem}[Pathwise positivity via Rouch\'e]
  \label{thm:main}
  Suppose $Q_j$ is a polynomial, $c$ is its contour centroid,
  $m_{\mathrm{R}}(Q_j,\partial D)>0$, and $Q_j(0)>\epsilon_j$.  Then
  $Q_j$ has no zeros in $D$, hence no zeros in $[0,1]$; since $Q_j(0)>0$
  we have $Q_j(t)>0$ for all $t\in[0,1]$.  If additionally
  $\max_{t\in[0,1]}\abs{g_j(t)-Q_j(t)}\leq\epsilon_j$, then
  $g_j(t)>0$ for all $t\in[0,1]$.
\end{theorem}

\begin{proof}
  $m_{\mathrm{R}}>0$ means $\abs{c}>\abs{Q_j(z)-c}$ for all $z\in\partial D$.
  By Rouch\'e's theorem, $Q_j$ and the constant $c\neq 0$ have the same
  number of zeros in $D$, namely zero.  Since $Q_j$ is continuous and zero-free
  on the connected set $[0,1]\subset D$, its sign is constant there and fixed
  by $Q_j(0)>\epsilon_j>0$.  The error bound then gives
  $g_j(t)\geq Q_j(t)-\epsilon_j>0$.
\end{proof}

%% ============================================================
\section{Counterexample to Naive Complexification}
\label{sec:counterexample}

A natural but incorrect approach is to extend the real path margin $g_j$
directly to a complex function by replacing $t$ with $z\in\mathbb{C}$.
For a neural network $f$, this requires complexifying the activations
(e.g.\ $\tanh(z)$ for $z\in\mathbb{C}$), which introduces singularities
and branch cuts not present on the real line.

\begin{proposition}[Spurious zeros]
  \label{prop:counterex}
  Let $g_j(t) = \tanh(2t-1)$ (a sigmoid-like margin, always positive for
  $t\in[0.5+\delta, 1]$ for any $\delta>0$).  The complex extension
  $G_j(z) = \tanh(2z-1)$ has poles at $z_k = \tfrac{1}{2}(1 + i\pi(k+\tfrac{1}{2}))$
  for $k\in\mathbb{Z}$.  For $w>\pi/4\approx 0.785$, these poles lie inside
  the rectangular domain $D$, making $G_j$ non-analytic there.
\end{proposition}

\begin{proof}
  $\tanh(z)$ has poles at $z = i\pi(k+\tfrac{1}{2})$, $k\in\mathbb{Z}$.
  Substituting $z\mapsto 2z-1$ shifts poles to $z_k = \tfrac{1}{2}(1+i\pi(k+\tfrac{1}{2}))$.
  For $k=0$: $\operatorname{Im}(z_0) = \pi/4\approx 0.785$.  Hence for $w>\pi/4$ the domain
  $D$ contains a pole, invalidating any Cauchy-based certificate.
\end{proof}

This shows why the surrogate approach is essential: we fit a polynomial $Q_j$
to \emph{samples} of $g_j$ on the real segment and extend only the polynomial
analytically, avoiding the nonlinearities of the network entirely.

%% ============================================================
\section{Experiments}
\label{sec:experiments}

\subsection{Exp P3-1: Synthetic polynomial classifiers}
\label{sec:exp_synthetic}

We generate random degree-$d$ polynomials $g_j$ with controlled minimum
value $\mu\in\{0.5, 0.2, 0.05\}$ and fit surrogates of matching degree.
We run 100 instances per cell and compare three methods: dense-sampling
ground truth, real-polynomial certificate, and Rouch\'e certificate.

Table~\ref{tab:synthetic} reports certification rates.
Figure~\ref{fig:cert_heatmap} shows heatmaps over the degree$\times$margin grid.

\begin{table}[t]
\centering
\caption{Synthetic polynomial classifiers: certification rates (\%) and
  mean margin by method.  100 instances per cell; surrogates of degree
  matching $g_j$.}
\label{tab:synthetic}
\begin{tabular}{ccrrrc}
\toprule
Degree & Min margin & Dense (\%) & Real poly (\%) & Rouch\'e (\%) & Mean real margin \\
\midrule
\multirow{3}{*}{2}
  & 0.50 & 100.0 & 100.0 & 46.0 & $+0.500$ \\
  & 0.20 & 100.0 & 100.0 &  7.0 & $+0.200$ \\
  & 0.05 & 100.0 & 100.0 &  1.0 & $+0.050$ \\
\midrule
\multirow{3}{*}{4}
  & 0.50 & 100.0 & 100.0 & 14.0 & $+0.500$ \\
  & 0.20 & 100.0 & 100.0 &  7.0 & $+0.200$ \\
  & 0.05 & 100.0 & 100.0 &  0.0 & $+0.050$ \\
\midrule
\multirow{3}{*}{8}
  & 0.50 & 100.0 & 100.0 &  3.0 & $+0.500$ \\
  & 0.20 & 100.0 & 100.0 &  0.0 & $+0.200$ \\
  & 0.05 & 100.0 & 100.0 &  0.0 & $+0.050$ \\
\bottomrule
\end{tabular}
\end{table}

\begin{figure}[t]
\centering
\includegraphics[width=\linewidth]{../../results/p3exp1/cert_rate_heatmap.pdf}
\caption{Certification rate heatmaps for synthetic polynomial classifiers.
  Left: dense ground truth.  Centre: real-polynomial certificate.
  Right: Rouch\'e certificate.  Colour scale 0--100\%.}
\label{fig:cert_heatmap}
\end{figure}

\subsection{Exp P3-2: Small MLP on two-moons}
\label{sec:exp_networks}

We train a two-hidden-layer MLP ($2\to 32\to 32\to 2$, Tanh activations,
Adam, 300 epochs) on the two-moons dataset.  For 200 same-class straight-line
path pairs from the test set, we fit Chebyshev surrogates of degree 4, 8,
and 16, and evaluate real-polynomial and Rouch\'e certificates.

Table~\ref{tab:networks} compares certification rates and mean approximation
error.  Figure~\ref{fig:twomoons} shows the decision boundary and per-method
cert rates.

\begin{table}[t]
\centering
\caption{Two-moons MLP: certification rates (\%) and mean surrogate
  approximation error by degree.  200 same-class paths; contour width $w=0.3$.}
\label{tab:networks}
\begin{tabular}{crrrc}
\toprule
Degree & Dense (\%) & Real poly (\%) & Rouch\'e (\%) & Mean $\epsilon_j$ \\
\midrule
 4 & 78.5 & 78.5 & 66.5 & $5.99\times10^{-4}$ \\
 8 & 78.5 & 78.5 & 67.0 & $2.64\times10^{-6}$ \\
16 & 78.5 & 78.5 & 34.5 & $6.96\times10^{-7}$ \\
\bottomrule
\end{tabular}
\end{table}

\subsection{Exp P3-3: Scaling limits}
\label{sec:exp_limits}

To stress-test the certificates we scale the interpolation path by factors
$s\in\{1, 2, 4, 8\}$, extending $\gamma(t)=x_a + st(x_b-x_a)$ beyond the
original endpoints.  Longer paths have higher Lipschitz variation and thus
larger surrogate approximation error.

Table~\ref{tab:scaling} and Figure~\ref{fig:twomoons} (right panel) show that
certification rates collapse rapidly with scale while approximation error grows.

\begin{table}[t]
\centering
\caption{Scaling limits: cert rates (\%) and mean approximation error vs
  path scale.  Degree 8 surrogate; 200 paths.}
\label{tab:scaling}
\begin{tabular}{crrc}
\toprule
Scale & Real poly (\%) & Rouch\'e (\%) & Mean $\epsilon_j$ \\
\midrule
1 & 78.5 & 63.0 & $3.82\times10^{-6}$ \\
2 & 65.5 & 36.5 & $3.90\times10^{-4}$ \\
4 & 53.0 & 16.0 & $5.05\times10^{-3}$ \\
8 & 48.0 &  6.0 & $2.89\times10^{-2}$ \\
\bottomrule
\end{tabular}
\end{table}

\begin{figure}[t]
\centering
\includegraphics[width=\linewidth]{../../results/p3exp2/two_moons_certs.pdf}
\caption{Left: certification rates vs surrogate degree for three methods.
  Right: certification rates and approximation error vs path scale (degree 8).}
\label{fig:twomoons}
\end{figure}

%% ============================================================
\section{Results}
\label{sec:results}

\paragraph{Synthetic polynomial classifiers (Table~\ref{tab:synthetic}).}
With surrogates of matching degree, the approximation error is numerically
zero, so the real-polynomial certificate achieves 100\% on all 9 cells.
The Rouch\'e certificate is dramatically more conservative: at degree~2 and
minimum margin 0.5 it certifies only 46\% of paths, and at degree~8 it
certifies 0--3\% regardless of margin.  The failure mode is purely geometric:
high-degree polynomials take large complex values off the real axis, so the
contour centroid $c$ satisfies $|c| \ll \max_{\partial D}|Q_j(z)-c|$ and the
Rouch\'e condition is never met.  The result is independent of margin: even
with minimum margin 0.5, Rouché certifies only 3\% at degree~8.

\paragraph{Two-moons MLP, method comparison (Table~\ref{tab:networks}).}
The dense ground-truth rate of 78.5\% is constant across all surrogate
degrees, reflecting the true fraction of same-class test pairs whose
straight-line interpolation stays within the predicted class.
Real-polynomial certificates match dense exactly at all degrees:
approximation errors are at most $6\times10^{-4}$, negligible relative to
path margins of order 0.1--1.  Rouché falls below real poly at every degree
(66.5\%, 67.0\%, 34.5\% vs.\ a constant 78.5\%), and collapses to 34.5\% at degree~16,
confirming that higher-degree surrogates amplify the complex excursion
problem observed in the synthetic setting.

\paragraph{Scaling limits (Table~\ref{tab:scaling}).}
With degree-8 surrogates and path scale increasing from 1$\times$ to
8$\times$, real-polynomial cert rate falls from 78.5\% to 48.0\% as longer
paths encounter decision boundaries.  The approximation error grows by four
orders of magnitude (from $3.8\times10^{-6}$ to $2.9\times10^{-2}$) due to
higher variation in $g_j$ over longer paths.  Rouch\'e collapses much faster:
from 63.0\% at scale~1 to 6.0\% at scale~8, because both the complex
excursion and the error growth compound.  At scale~8 the Rouch\'e certificate
is essentially vacuous while the real-polynomial check still certifies
nearly half of all paths.

\paragraph{Summary.}
Across all settings, the real-polynomial positivity check equals or exceeds
the Rouch\'e certificate in certification rate, with no exception.  The
Rouché certificate's complex-domain geometry provides no practical advantage
and introduces a second failure mode (contour excursion) absent from the
real check.

%% ============================================================
\section{Discussion and Practical Guidelines}
\label{sec:discussion}

\subsection{When Rouch\'e certificates are informative}

The real-polynomial certificate strictly dominates Rouché in every
experimental setting.  On synthetic classifiers with exact surrogates
(approximation error $\approx 0$), real-polynomial achieves 100\% while
Rouché achieves 0--46\%; the gap widens with polynomial degree (3\% at
degree~8 even with minimum margin~0.5).  The Rouché certificate adds
computational cost (complex contour evaluation) and introduces a second
failure mode absent from the real check.

For Rouch\'e to offer a strict advantage over the real-polynomial check, the
polynomial $Q_j$ would need to have its minimum on the real segment obscured by
local oscillation while remaining robustly positive everywhere in the complex
domain $D$.  This configuration does not arise naturally for smooth margin
functions, because polynomials that are large on $\partial D$ tend also to be
large on the real interior.

\subsection{Why Rouch\'e fails on neural paths}

For neural network margin functions, the fundamental obstacle is the
approximation error $\epsilon_j$.  A degree-$d$ surrogate on a path of length
$L$ (measured in input space) achieves error
$\epsilon_j = O(L^d \cdot \omega(g_j, L/d))$
where $\omega$ is the modulus of continuity.  As $L$ grows, both the variation
of $g_j$ and the required degree increase, but the margin $\min_t g_j(t)$
typically decreases near class boundaries.  The certification condition
$\min Q_j > \epsilon_j$ becomes impossible to satisfy.

\subsection{Practical guidelines}

\begin{enumerate}
  \item \textbf{Short interpolation paths only.}  Certification is feasible
    only for short same-class interpolations with margin $\gg \epsilon_j$.
  \item \textbf{Real-polynomial check is sufficient.}  The Rouch\'e certificate
    adds no power beyond $\min_{[0,1]} Q_j > \epsilon_j$ for typical smooth
    margins.  Prefer the simpler check.
  \item \textbf{Rigorise the error bound.}  Our $\epsilon_j$ estimate is
    numerical; a verified bound via Bernstein ellipse analysis or interval
    arithmetic is needed for rigorous guarantees.
  \item \textbf{Consider Lipschitz-based methods.}  For neural networks with
    known Lipschitz constants, direct Lipschitz certificates
    $\min_t g_j(t) > L_j \cdot \norm{x_b - x_a}$ are tighter and require no
    surrogate fitting.
\end{enumerate}

%% ============================================================
\section{Conclusion}
\label{sec:conclusion}

We have built and evaluated a complete pipeline for Rouch\'e-type path
certificates: Chebyshev surrogate fitting, approximation error estimation,
constant-reference Rouch\'e margin, and pathwise positivity theorem.  The
theoretical framework is sound---Theorem~\ref{thm:main} provides a rigorous
connection between complex zero-counting and real path positivity, and
Proposition~\ref{prop:counterex} explains why naive complexification fails.

Empirically, the real-polynomial certificate strictly dominates Rouché
in every setting tested: on synthetic classifiers with exact surrogates
real-polynomial achieves 100\% while Rouché achieves 0--46\%; on
two-moons paths Rouché lags 12--44 percentage points behind real-polynomial
across all degrees and scales.  On realistic neural paths beyond short
interpolation scales, both certificates become vacuous as surrogate
approximation error consumes the available margin, with Rouché collapsing
faster due to the compounding complex-excursion failure mode.

This negative result clarifies the boundary between elegant complex analysis
and practical neural network verification: Rouch\'e zero-counting is
well-suited to settings where the analytic structure is intrinsic (as in
Papers~1 and~2, where the contour is the unit circle and the matrix
characteristic function is already analytic), but not to settings where
analyticity must be imposed through a surrogate on an arbitrary real path.

\end{document}
